\section{模型与产品：从裸模型到 Cloud/OS/容器}

上一章讨论 Online Learning，本章转向商业和产品。广密认为裸模型发布的时代可能逐渐弱化，壁垒会更多来自 Cloud、OS、生态和工作流容器。原因是模型能力会被追赶、开源冲击和 API 化，但用户每天在哪个环境里完成任务，哪个系统拥有状态、数据、权限和支付关系，才会形成长期壁垒。

\begin{figure}[H]
\centering
\includegraphics[width=0.92\textwidth]{model-product-moat.png}
\caption{模型与产品的壁垒：裸模型发布时代弱化，Cloud/OS/生态成为防线。自制概念图，依据 01:02:45--01:15:11 对谈内容整理。}
\end{figure}

\begin{knowledgebox}{读图：模型强不等于产品强}
左侧 Model 强调能力上限和成本效率，右侧 Product/OS 强调入口、工作流、生态、计费和数据回流。读这张图时要抓住一点：模型是发动机，产品是道路、仪表盘、收费系统和驾驶体验。
\end{knowledgebox}

\subsection{20 美元定价与 SaaS 复制问题}

本节看定价。访谈追问为什么很多 AI 产品定价接近 20 美元，是否只是复制 SaaS。广密的隐含问题是：AI 产品的边际成本与传统 SaaS 不同。一个复杂 Agent 任务可能消耗大量 token 和工具调用，如果按低价包月无节制使用，就会出现毛利压力。因此未来定价可能更接近用量、任务价值、企业席位和算力消耗的混合模型。

\begin{warningbox}{定价不是 UI 问题，是计算经济问题}
如果产品把高成本 Agent 任务包装成低价无限用，短期可能获得增长，长期会被推理成本和用户滥用压垮。AI 产品经理必须理解 token、缓存、工具调用和任务价值。
\end{warningbox}

\subsection{模型会吃掉产品吗}

本节讨论产品公司最担心的问题。广密把本质问题概括为 feature system vs learning system。传统产品的 feature 可以被模型公司复制，但如果产品公司沉淀的是环境、数据、工作流、用户状态、权限体系和反馈闭环，就不容易被一个新模型版本直接吃掉。反过来，如果产品只是给模型套壳，没有独特环境和数据，确实容易被底座模型吞掉。

\begin{importantbox}{应用壁垒的判断标准}
看一个 AI 应用是否有壁垒，不要只看界面和 prompt，而要看它是否拥有：真实任务入口、用户长期状态、可验证反馈、专有工作流、成本控制、数据回流和替换成本。
\end{importantbox}

\begin{figure}[H]
\centering
\includegraphics[width=0.92\textwidth]{fire-thieves.png}
\caption{模型盗火者：Perplexity、Cursor、Manus 把前沿模型能力搬进具体工作流。自制概念图，依据 01:02:45--01:15:11 对谈内容整理。}
\end{figure}

\begin{knowledgebox}{读图：盗火者的价值在容器}
Perplexity 把模型能力放进搜索/研究，Cursor 放进编程环境，Manus 放进通用执行。它们不是自己训练最强底座模型，而是把模型能力放进用户愿意付费、愿意反复使用、能产生反馈的容器。
\end{knowledgebox}

\subsection{投资人应该看什么}

本节把产品问题转成投资判断。广密提到理想组合和公司权重，但本笔记不把它当投资建议。更重要的是分析框架：底座模型公司看能力上限、成本曲线和组织稳定性；应用公司看是否抓住模型能力外溢窗口、是否沉淀工作流、是否有真实收入和可控成本；基础设施公司看是否控制环境、推理、训练、数据和部署关键环节。

\begin{knowledgebox}{课堂提示：科技投资不是押热词}
访谈最后强调“创造”而不是“混圈”。在 AI 投资里，这意味着要理解技术路线和价值链：谁提升智能，谁承接智能，谁提供环境，谁掌握反馈，谁只是叙事套利。
\end{knowledgebox}

\subsection{本章小结}

模型与产品的关系不是谁吃掉谁这么简单。强模型会压缩薄产品空间，但也会创造新容器机会。产品壁垒来自环境、状态、反馈、成本和用户任务，而不是简单套一个聊天框。

